% GENERATED FILE. Edit canonical lessons, then run npm run generate:books.
% canonical-source-hash: fnv1a64:cf910eae8ad4f218
% canonical-lessons: IT-C457-cassiere, IT-C457-spiccioli, IT-C457-camerino, IT-C457-svendita, IT-C457-mercanteggiare

\chapter{Alla cassa}
\label{ch:it-alla-cassa}
\hlchaptermodality{457}

\begin{chapteropening}
\textbf{By the end of this chapter:} \emph{I can pay and try things on: the cashier, small change, the fitting room, a clearance sale and haggling.}

Complete the last lesson of chapter 457: I can pay and try things on: the cashier, small change, the fitting room, a clearance sale and haggling.
\end{chapteropening}

\section[il cassiere]{il cassiere — the cashier}

\label{lesson:IT-C457-cassiere}



\begin{quote}
Before the new one: say the Italian for the optician, then the Italian for to remove stains from.
\end{quote}

\subsection*{You'll want to know: il cassiere}
\textbf{il cassiere} — \textquotedblleft{}the cashier\textquotedblright{}.

\begin{quote}
Il cassiere mi chiede la tessera.
\end{quote}

\begin{quote}
\emph{The cashier asks me for my loyalty card.}
\end{quote}

\emph{La cassa} — the till; a woman is \emph{la cassiera}.

\begin{cousinweb}
\emph{cassiere} is built on \emph{cassa} 'till, box', which is inherited from Latin \emph{capsa} 'box' (EtymDB).
\end{cousinweb}

\begin{grammarlens}[title={What you've built}]
One more word for telling, explaining and describing at B1.
\end{grammarlens}

\subsection*{Your turn}
\begin{itemize}
\raggedright
  \item \emph{Say it:} \emph{il cassiere}
  \item \emph{Say it:} \emph{il cassiere}, once more
  \item \emph{Say it:} a sentence of your own with \emph{il cassiere}
  \item \emph{Recall:} say the Italian for the optician, then the Italian for to remove stains from, then say \emph{il cassiere} again
\end{itemize}

\begin{tcolorbox}[breakable,colback=teal!4,colframe=teal!35!black,title={Before you move on}]
What does il cassiere mean? (\textquotedblleft{}The cashier\textquotedblright{}.) Say it once more.
\end{tcolorbox}
\section[gli spiccioli]{gli spiccioli — the small change, the coins}

\label{lesson:IT-C457-spiccioli}



\begin{quote}
Before the new one: say the Italian for to remove stains from, then the Italian for the cashier.
\end{quote}

\subsection*{You'll want to know: gli spiccioli}
\textbf{gli spiccioli} — \textquotedblleft{}the small change, the coins\textquotedblright{}.

\begin{quote}
Non ho spiccioli per il caffè.
\end{quote}

\begin{quote}
\emph{I have no change for the coffee.}
\end{quote}

The singular \emph{lo spicciolo} is rare.

\begin{cousinweb}
\emph{spiccioli} is the plural of \emph{spicciolo} 'small coin'. The origin is not settled; it is often linked to \emph{spicciare} 'to sort out'.
\end{cousinweb}

\begin{grammarlens}[title={What you've built}]
One more word for telling, explaining and describing at B1.
\end{grammarlens}

\subsection*{Your turn}
\begin{itemize}
\raggedright
  \item \emph{Say it:} \emph{gli spiccioli}
  \item \emph{Say it:} \emph{gli spiccioli}, once more
  \item \emph{Say it:} a sentence of your own with \emph{gli spiccioli}
  \item \emph{Recall:} say the Italian for to remove stains from, then the Italian for the cashier, then say \emph{gli spiccioli} again
\end{itemize}

\begin{tcolorbox}[breakable,colback=teal!4,colframe=teal!35!black,title={Before you move on}]
What does gli spiccioli mean? (\textquotedblleft{}The small change, the coins\textquotedblright{}.) Say it once more.
\end{tcolorbox}
\section[il camerino]{il camerino — the fitting room; the dressing room}

\label{lesson:IT-C457-camerino}



\begin{quote}
Before the new one: say the Italian for the cashier, then the Italian for the small change.
\end{quote}

\subsection*{You'll want to know: il camerino}
\textbf{il camerino} — \textquotedblleft{}the fitting room; the dressing room\textquotedblright{}.

\begin{quote}
Posso provarlo in camerino?
\end{quote}

\begin{quote}
\emph{Can I try it on in the fitting room?}
\end{quote}

In a theatre, \emph{il camerino} is the actors' dressing room.

\begin{cousinweb}
\emph{camerino} is a little \emph{camera} 'room', which is inherited from Latin \emph{camera} 'vault, arched roof' (EtymDB).
\end{cousinweb}

\begin{grammarlens}[title={What you've built}]
One more word for telling, explaining and describing at B1.
\end{grammarlens}

\subsection*{Your turn}
\begin{itemize}
\raggedright
  \item \emph{Say it:} \emph{il camerino}
  \item \emph{Say it:} \emph{il camerino}, once more
  \item \emph{Say it:} a sentence of your own with \emph{il camerino}
  \item \emph{Recall:} say the Italian for the cashier, then the Italian for the small change, then say \emph{il camerino} again
\end{itemize}

\begin{tcolorbox}[breakable,colback=teal!4,colframe=teal!35!black,title={Before you move on}]
What does il camerino mean? (\textquotedblleft{}The fitting room; the dressing room\textquotedblright{}.) Say it once more.
\end{tcolorbox}
\section[la svendita]{la svendita — the clearance sale}

\label{lesson:IT-C457-svendita}



\begin{quote}
Before the new one: say the Italian for the small change, then the Italian for the fitting room.
\end{quote}

\subsection*{You'll want to know: la svendita}
\textbf{la svendita} — \textquotedblleft{}the clearance sale\textquotedblright{}.

\begin{quote}
C'è una svendita: tutto a metà prezzo.
\end{quote}

\begin{quote}
\emph{There is a clearance sale: everything half price.}
\end{quote}

\emph{I saldi} are the seasonal sales.

\begin{cousinweb}
\emph{svendita} is built from \emph{s-} and \emph{vendita} 'sale', from \emph{vendere} 'to sell', which is inherited from Latin \emph{vēndere} (EtymDB).
\end{cousinweb}

\begin{grammarlens}[title={What you've built}]
One more word for telling, explaining and describing at B1.
\end{grammarlens}

\subsection*{Your turn}
\begin{itemize}
\raggedright
  \item \emph{Say it:} \emph{la svendita}
  \item \emph{Say it:} \emph{la svendita}, once more
  \item \emph{Say it:} a sentence of your own with \emph{la svendita}
  \item \emph{Recall:} say the Italian for the small change, then the Italian for the fitting room, then say \emph{la svendita} again
\end{itemize}

\begin{tcolorbox}[breakable,colback=teal!4,colframe=teal!35!black,title={Before you move on}]
What does la svendita mean? (\textquotedblleft{}The clearance sale\textquotedblright{}.) Say it once more.
\end{tcolorbox}
\section[mercanteggiare]{mercanteggiare — to haggle, to bargain}

\label{lesson:IT-C457-mercanteggiare}



\begin{quote}
Before the new one: say the Italian for the fitting room, then the Italian for the clearance sale.
\end{quote}

\subsection*{You'll want to know: mercanteggiare}
\textbf{mercanteggiare} — \textquotedblleft{}to haggle, to bargain\textquotedblright{}.

\begin{quote}
Al mercato si può mercanteggiare un po'.
\end{quote}

\begin{quote}
\emph{At the market you can haggle a bit.}
\end{quote}

More common in speech: \emph{tirare sul prezzo}.

\begin{cousinweb}
\emph{mercanteggiare} is built on \emph{mercante} 'merchant', which is inherited from a form of Latin \emph{mercārī} 'to trade', from \emph{merx} 'goods' (Lewis \& Short).
\end{cousinweb}

\begin{grammarlens}[title={What you've built}]
One more word for telling, explaining and describing at B1.
\end{grammarlens}

\subsection*{Your turn}
\begin{itemize}
\raggedright
  \item \emph{Say it:} \emph{mercanteggiare}
  \item \emph{Say it:} \emph{mercanteggiare}, once more
  \item \emph{Say it:} a sentence of your own with \emph{mercanteggiare}
  \item \emph{Recall:} say the Italian for the fitting room, then the Italian for the clearance sale, then say \emph{mercanteggiare} again
\end{itemize}

\begin{tcolorbox}[breakable,colback=teal!4,colframe=teal!35!black,title={Before you move on}]
What does mercanteggiare mean? (\textquotedblleft{}To haggle, to bargain\textquotedblright{}.) Say it once more.
\end{tcolorbox}
